\chapter{Hashtables: Closed addressing}\label{chap:hashtables:closed}

Closed addressing (also known as chaining) handles collisions by storing all elements that hash to the same bucket in a secondary data structure, typically a linked list.

\section{Design rationale}

Our implementation uses a dynamic array of buckets, where each bucket is a pointer to the head of a linked list. We reuse the generic list implementation described in Chapter~\ref{chap:linear}.

Advantages of this approach include:
\begin{itemize}
    \item \textbf{Simplicity:} The logic for handling collisions is straightforward.
    \item \textbf{Graceful degradation:} Performance degrades linearly with the load factor, even if it exceeds 1.0.
    \item \textbf{Deletion:} Deleting an element is as simple as removing it from the linked list.
\end{itemize}

\section{Example variant: Move-to-Front}

To improve performance for temporal locality (where recently accessed elements are likely to be accessed again), we implemented a Move-to-Front heuristic. Upon a successful \texttt{get} operation, the found element is moved to the front of its collision list. This reduces the search time for subsequent lookups of the same key.
